\DocumentMetadata{
  pdfversion=2.0,
  pdfstandard=ua-2,
  lang=en-US,
  tagging=on,
  tagging-setup={math/setup=mathml-SE,tabsorder=structure}
}
\documentclass[11pt,letterpaper]{article}
\usepackage[margin=0.68in,includefoot]{geometry}
\usepackage{newcomputermodern}
\usepackage{amsmath,amscd,mathtools}
\usepackage{tikz,adjustbox}
\providecommand{\square}{\mdlgwhtsquare}
\usepackage[hidelinks]{hyperref}
\hypersetup{
  pdftitle={Numerical Analysis PhD Exam, January 2016, Part 2},
  pdfauthor={University of Florida Department of Mathematics},
  pdfsubject={Graduate mathematics examination},
  pdfdisplaydoctitle=true,
  bookmarksopen=true
}
\setcounter{secnumdepth}{0}
\setlength{\parindent}{0pt}
\setlength{\parskip}{0.28em}
\setlength{\emergencystretch}{3em}
\setlength{\leftmargini}{2.15em}
\setlength{\leftmarginii}{2.65em}
\setlength{\leftmarginiii}{3em}
\renewcommand{\labelenumi}{\textbf{\theenumi.}}
\pagestyle{empty}
\raggedbottom
\allowdisplaybreaks
\newenvironment{examproblems}[1]{%
  \begin{enumerate}%
  \setcounter{enumi}{#1}%
  \setlength{\topsep}{0.35em}%
  \setlength{\partopsep}{0pt}%
  \setlength{\itemsep}{0.48em}%
  \setlength{\parsep}{0.18em}%
}{\end{enumerate}}
\newenvironment{examsubparts}{%
  \begin{enumerate}%
  \setlength{\topsep}{0.18em}%
  \setlength{\partopsep}{0pt}%
  \setlength{\itemsep}{0.16em}%
  \setlength{\parsep}{0.1em}%
}{\end{enumerate}}
\newenvironment{examinstructions}{%
  \par\begingroup\itshape
}{%
  \par\endgroup\vspace{0.18em}
}
\makeatletter
\renewcommand\section{\@startsection{section}{1}{0pt}%
  {-1.25ex plus -.25ex minus -.1ex}%
  {0.55ex plus .1ex}%
  {\normalfont\large\bfseries}}
\renewcommand{\@maketitle}{%
  \newpage\null\vskip -1em%
  {\footnotesize\bfseries
    \href{https://www.ufl.edu/}{University of Florida}, \href{https://math.ufl.edu/}{Department of Mathematics}\par}%
  \vskip -0.5em%
  \tagstructbegin{tag=Title}%
    {\LARGE\bfseries\href{https://gma.math.ufl.edu/exams/numerical-analysis-phd/}{Numerical Analysis PhD Exam}, January 2016, Part 2\par}%
  \tagstructend%
  \vskip 0.25em%
  \tagmcbegin{artifact}%
    \hrule height 1pt%
  \tagmcend%
  \vskip 1em%
}
\makeatother
\title{Numerical Analysis PhD Exam, January 2016, Part 2}
\author{}
\date{}
\begin{document}
\maketitle
\thispagestyle{empty}
\begin{examinstructions}
You must show all your work as neatly and clearly as possible and indicate the final answer clearly.
\end{examinstructions}
\begin{examproblems}{0}
\item
A fourth degree polynomial \(P(x)\) satisfies \(\Delta^4P(0)=24\), \(\Delta^3P(0)=6\), and \(\Delta^2P(0)=0\), where \(\Delta P(x)=P(x+1)-P(x)\). Compute \(\Delta^2P(10)\).
\item
Consider the fixed point iteration 
\[
x_{n+1}=\phi(x_n),\qquad n=0,1,2,\ldots
\]
 where 
\[
\phi(x)=Ax+Bx^2+Cx^3.
\]
 Given a positive number~\(\alpha\), determine the constants \(A,B,C\) such that the iteration converges locally to \(1/\alpha\) with order \(p=3\).
\item
This problem has the following parts:
\begin{examsubparts}
\item[\textnormal{(a)}]
Find~\(\alpha\), \(\beta\) and~\(\gamma\) so that the quadrature formula has a maximum degree of precision. What is the exact degree of precision of this quadrature formula? 
\[
\int_0^2 f(x)\,dx\approx \alpha(f(0)+f(2))+\beta(f'(0)-f'(2))+\gamma(f''(0)-f''(2)).
\]
\item[\textnormal{(b)}]
Suggest your own quadrature formula for the integral 
\[
\int_0^2 f(x)\,dx.
\]
 Your formula should have expected degree of precision at least four (you don’t have to compute to demonstrate that).
\end{examsubparts}
\item
For the function 
\[
f(x)=\sqrt[m]{|x|},\qquad m\ne\frac12,
\]
 in the interval \([-1,1]\), find the polynomial of minimax approximation of degree 2. What is the minimax error? What happens if \(m=\frac12\)?
\item
Prove that Gaussian quadrature formula (for any~\(n\)) in the interval \([-1,1]\) and with weight \(w(x)=1\) is exact for all odd functions.
\end{examproblems}
\end{document}
